\documentclass[11pt]{article}

\usepackage[margin=1in]{geometry}
\usepackage{amsmath, amssymb}
\usepackage{tikz}
\usetikzlibrary{arrows.meta, positioning, fit, backgrounds, calc}
\usepackage{float}
\usepackage{tcolorbox}
\tcbuselibrary{skins, breakable}
\usepackage{xcolor}
\usepackage{enumitem}
\usepackage{hyperref}
\hypersetup{colorlinks=true, linkcolor=blue!55!black, urlcolor=blue!55!black}

\definecolor{accent}{RGB}{30,90,160}
\definecolor{accent2}{RGB}{170,60,30}
\definecolor{soft}{RGB}{238,243,250}
\definecolor{soft2}{RGB}{250,243,236}

\newtcolorbox{keybox}[1]{
  colback=soft, colframe=accent, fonttitle=\bfseries,
  title=#1, breakable, boxrule=0.6pt, arc=2pt, left=6pt, right=6pt}

\newtcolorbox{watchbox}[1]{
  colback=soft2, colframe=accent2, fonttitle=\bfseries,
  title=#1, breakable, boxrule=0.6pt, arc=2pt, left=6pt, right=6pt}

\newcommand{\dd}[2]{\frac{\partial #1}{\partial #2}}
\newcommand{\ddt}[1]{\frac{d#1}{dt}}

\title{\textbf{Mann \S4.2--4.3 --- Noether's Theorem \&\\
Spacetime Currents}\\[2pt]
\large A Plain--Language Walkthrough}
\author{Reading companion --- Mann, \emph{An Introduction to Particle Physics
and the Standard Model}, Ch.~4}
\date{2026--06--16}

\begin{document}
\maketitle

\noindent\textbf{Where we are.} In \S4.1 we learned the \emph{action principle}:
the true path of a system extremizes the action $S[q]=\int dt\,L[q,\dot q]$,
which gives the Euler--Lagrange equations and defines the canonical momentum
$p_i=\partial L/\partial\dot q_i$. We also saw that if a transformation leaves
$S$ unchanged (\emph{invariance}), the equations of motion keep their form
(\emph{covariance}). Now \S4.2 turns that into the single most useful theorem
in the chapter, and \S4.3 reads off the three classic conservation laws from it.

\bigskip
\hrule
\bigskip

\section*{1. \S4.2 Noether's Theorem --- the one-sentence version}

\begin{keybox}{The whole idea}
\textbf{Every continuous symmetry of the action comes with a quantity that does
not change in time.} Symmetry in $\Rightarrow$ conservation law out. (Emmy
Noether, 1917.)
\end{keybox}

\noindent Why should that be true? The proof is short once you spot the trick.

\subsection*{1.1 The trick: allow a ``sloppy'' variation}

Back in \S4.1 we varied the path but imposed two tidy conditions: the
equations of motion held, \emph{and} the endpoints were pinned ($\delta q=0$ at
the ends). Noether's proof does the opposite --- it allows a \emph{fully
arbitrary} variation, with \emph{neither} condition imposed. Carrying the
algebra of \S4.1 one step further (an integration by parts, but now \emph{keeping}
the boundary term) gives the master identity
\begin{equation}
\delta_{q_i} L
= \underbrace{\left[\dd{L}{q_i}-\ddt{}\dd{L}{\dot q_i}\right]}_{\text{Euler--Lagrange bracket}}\delta q_i
\;+\; \ddt{}\big[p_i\,\delta q_i\big].
\label{eq:master}
\end{equation}

\begin{keybox}{How to read eq.~\eqref{eq:master}}
The change in the Lagrangian splits into two pieces:
\begin{itemize}[nosep]
  \item the \textbf{Euler--Lagrange bracket} --- this is \emph{zero whenever the
  real equations of motion hold};
  \item a \textbf{total time derivative} $\dfrac{d}{dt}\big[p_i\,\delta q_i\big]$
  --- a ``boundary'' piece that is usually \emph{not} a pure $d/dt$ of anything
  nice.
\end{itemize}
\end{keybox}

\subsection*{1.2 Now feed in a symmetry}

Describe a small symmetry by parameters $\alpha=\{\alpha_1,\dots,\alpha_m\}$,
with $\alpha=0$ meaning ``do nothing'' (the identity). The induced change in the
coordinates is
\begin{equation}
\delta_{\alpha_j} q_i = \left.\dd{q_i'}{\alpha_j}\right|_{\alpha=0}\delta\alpha_j .
\end{equation}

Here is the pivotal fact that makes a symmetry \emph{special}:

\begin{keybox}{Symmetry $\Leftrightarrow$ $\delta L$ is a total time derivative}
If the transformation is a genuine symmetry, then $S$ is unchanged, so
$\delta_\alpha S=\int dt\,\delta_\alpha L=0$. The most general way an integral can
vanish for arbitrary limits is for the integrand itself to be a total
derivative:
\begin{equation}
\delta_{\alpha_k} L = \ddt{}\,G\!\left(q,\dd{q'}{\alpha}\right),
\end{equation}
where $G$ is some function (vanishing at the endpoints) that you must work out
\textbf{case by case}. \emph{For a symmetry $\delta_\alpha L$ is always a total
$d/dt$; for a random variation it is not.} That single difference is the engine
of the whole theorem.
\end{keybox}

\subsection*{1.3 Build the conserved current}

Combine the two total-derivative facts. Define the \textbf{Noether current}
\begin{equation}
\boxed{\,J_k = p_i\,\dd{q_i'}{\alpha_k} - G\,}\qquad\text{(one $J_k$ per symmetry parameter $\alpha_k$).}
\label{eq:current}
\end{equation}
Differentiate it in time. Using eq.~\eqref{eq:master} \emph{and} imposing the
Euler--Lagrange equations (now we \emph{do} sit on a real solution), the
Euler--Lagrange bracket dies, and the leftover is exactly $\frac{d}{dt}G$, which
cancels the $-G$ term:
\begin{equation}
\ddt{J_k}
= \ddt{}\!\left(p_i\dd{q_i'}{\alpha_k}\right)-\ddt{G}
= \underbrace{\left[\dd{L}{q_i}-\ddt{}\dd{L}{\dot q_i}\right]}_{=\,0\ \text{(EOM)}}\dd{q_i'}{\alpha_k}
\;+\;\ddt{G}-\ddt{G}
= 0 .
\end{equation}

\begin{keybox}{Noether's theorem (the payoff)}
\[
\ddt{J_k}=0 \qquad\Longrightarrow\qquad J_k \text{ is \textbf{conserved}.}
\]
What $J_k$ \emph{physically means} depends entirely on \emph{which} symmetry you
plugged in. Section 4.3 plugs in the three spacetime symmetries and reads off
the three famous conservation laws.
\end{keybox}

\begin{watchbox}{Two easy-to-miss points}
\begin{itemize}[nosep]
  \item \textbf{Two different variations are used.} The master identity
  \eqref{eq:master} comes from a \emph{sloppy} variation (EOM not assumed). The
  final $dJ_k/dt=0$ is evaluated \emph{on} a real solution (EOM assumed). Both
  steps are needed.
  \item \textbf{$G$ is not optional decoration.} For translations/rotations it
  turns out $G=0$, but for time translation $G=L$. Getting $G$ right is what
  makes $J_k$ truly constant.
\end{itemize}
\end{watchbox}

\bigskip
\hrule
\bigskip

\section*{2. \S4.3 Spacetime Symmetries and Their Noether Currents}

Now we just turn the crank of eq.~\eqref{eq:current} three times. Each row of
the table below is one symmetry; the punchline conservation laws are the three
pillars of mechanics.

\subsection*{2.1 Spatial translations $\Rightarrow$ linear momentum}

\textbf{Physical statement:} there is no special point in space; shifting the
origin changes nothing. So the action is invariant under $\vec q\to\vec q+\vec\alpha$:
\begin{equation}
L[\,\vec q+\vec\alpha,\dot{\vec q}\,]=L[\,\vec q,\dot{\vec q}\,].
\end{equation}
Three parameters $\{\alpha_1,\alpha_2,\alpha_3\}$ (the shift distances). Because
$L$ does not change, $\partial L/\partial\alpha_i|_0=0$, so we may take $G=0$.
Also $\partial q_i'/\partial\alpha_k=\delta_{ik}$. Plug into \eqref{eq:current}:
\begin{equation}
J_k^{\text{sp-trans}} = p_i\,\delta_{ik}-0 = p_k
\qquad\Longrightarrow\qquad \ddt{\vec p}=0 .
\end{equation}

\begin{keybox}{Result 1}
\[
\textbf{Space-translation invariance}\ \Longrightarrow\ \textbf{conservation of linear momentum }\vec p .
\]
\end{keybox}

\subsection*{2.2 Rotations $\Rightarrow$ angular momentum}

\textbf{Physical statement:} no special orientation in space. Rotations form the
group $SO(3)$; for a small rotation (from Ch.~3),
\begin{equation}
q_i' = \big[e^{\,i\vec\alpha\cdot\vec J}\big]_i{}^{j} q_j
\;\simeq\; q_i-(\vec\alpha\times\vec q)_i ,
\qquad\text{so}\qquad
\dd{q_i'}{\alpha_k}=\epsilon_{ikj}\,q_j .
\end{equation}
A rotationally invariant $L$ can only depend on the \emph{magnitude} $|\vec q|$
(rotations preserve lengths), so again $\partial L/\partial\alpha_i|_0=0$ and
$G=0$. Then
\begin{equation}
J_k^{\text{rot}} = p_i\,\epsilon_{ikj}q_j = (\vec q\times\vec p)_k
\qquad\Longrightarrow\qquad \ddt{}(\vec q\times\vec p)=0 .
\end{equation}

\begin{keybox}{Result 2}
\[
\textbf{Rotational invariance}\ \Longrightarrow\ \textbf{conservation of angular momentum }\vec q\times\vec p .
\]
\end{keybox}

\subsection*{2.3 Time translations $\Rightarrow$ energy}

\textbf{Physical statement:} the laws of physics are the same today as
yesterday; shifting the clock's zero changes nothing. Under $t'=t+\delta\alpha$,
\begin{equation}
\left.\dd{q_i'}{\alpha}\right|_{0}=\dot q_i .
\end{equation}
Just \emph{one} parameter $\alpha$ (how far the time origin moves). This case is
different: $G\neq 0$. Because shifting time shifts the whole Lagrangian,
\begin{equation}
\delta_\alpha L = \ddt{L}\,\delta\alpha+\cdots
\;\Longrightarrow\; \left.\dd{L}{\alpha}\right|_0=\ddt{L}
\;\Longrightarrow\; G = L .
\end{equation}
Plugging $\partial q_i'/\partial\alpha=\dot q_i$ and $G=L$ into \eqref{eq:current}:
\begin{equation}
J^{\text{t-trans}} = p_i\,\dot q_i - L \equiv H
\qquad\Longrightarrow\qquad \ddt{H}=0 .
\end{equation}
The combination $H=p_i\dot q_i-L$ is exactly the \textbf{Hamiltonian} (the
energy).

\begin{keybox}{Result 3}
\[
\textbf{Time-translation invariance}\ \Longrightarrow\ \textbf{conservation of energy }H .
\]
\end{keybox}

\begin{watchbox}{Why energy conservation means a closed system}
Time translation is a symmetry exactly when $L$ has \emph{no explicit}
$t$-dependence. If $L$ depended on $t$, something external could be pumping the
system at a later time, and the action would shift when you move the clock.
``No explicit $t$'' is the definition of a \textbf{closed system} --- so energy
conservation holds for any closed system.
\end{watchbox}

\bigskip
\hrule
\bigskip

\section*{3. One-page recap}

\begin{keybox}{The \S4.2--4.3 machine}
\renewcommand{\arraystretch}{1.4}
\begin{center}
\begin{tabular}{@{}llll@{}}
\textbf{Symmetry} & $\dfrac{\partial q_i'}{\partial\alpha_k}$ & $G$ & \textbf{Conserved $J_k$}\\[4pt]
\hline
Space translation & $\delta_{ik}$ & $0$ & $p_k$ \;(linear momentum)\\
Rotation $SO(3)$   & $\epsilon_{ikj}q_j$ & $0$ & $(\vec q\times\vec p)_k$ \;(angular momentum)\\
Time translation   & $\dot q_i$ & $L$ & $H=p_i\dot q_i-L$ \;(energy)\\
\hline
\end{tabular}
\end{center}
\smallskip
\textbf{Recipe:} (1) write the symmetry and its parameters $\alpha$; (2) compute
$\partial q'/\partial\alpha$; (3) find $G$ from $\delta_\alpha L=\frac{d}{dt}G$;
(4) assemble $J_k=p_i\,\partial q_i'/\partial\alpha_k-G$; (5) it's automatically
conserved on solutions.
\end{keybox}

\end{document}
